% Colours of the presentation template, defined here as well so that the TikZ diagrams also compile in the cheat sheet.
\definecolor{nyupurple}{RGB}{90, 35, 140}
\definecolor{lightnyupurple}{RGB}{159, 120, 184}

% LSTM diagram, built in stages: #1 = 1 shows the cell-state highway only, 2 adds the forget and input gates
% with the candidate, 3 adds the output gate. All gates read the same pair (h_{t-1}, x_t), drawn once as a bus.
\providecommand{\lstmdiagram}[1]{%
\adjustbox{max width=\linewidth, max totalheight=0.42\textheight}{%
\begin{tikzpicture}[>=latex, thick, every node/.style={font=\small}, st/.style={draw, circle, minimum size=9mm, inner sep=1pt}, gt/.style={draw, rounded corners, fill=orange!25, minimum width=11mm, minimum height=7mm, align=center, inner sep=2pt}, op/.style={draw, circle, minimum size=6mm, inner sep=0pt, fill=white}]
\node[st] (cp) at (0,3) {$c_{t-1}$};
\node[op] (fm) at (3,3) {$\times$};
\node[op] (ad) at (6,3) {$+$};
\node[st] (ct) at (11,3) {$c_t$};
\node[font=\scriptsize, gray, above] at (cp.north) {old memory};
\node[font=\scriptsize, gray, above] at (ct.north) {new memory};
\draw[->, line width=1.6pt, lightnyupurple!70!black] (cp) -- (fm);
\draw[->, line width=1.6pt, lightnyupurple!70!black] (fm) -- (ad);
\draw[->, line width=1.6pt, lightnyupurple!70!black] (ad) -- (ct);
\ifnum#1=1
\node[font=\scriptsize, above=3pt] at (fm.north) {scale down};
\node[font=\scriptsize, above=3pt] at (ad.north) {add};
\fi
\ifnum#1>1
\node[font=\scriptsize, above=3pt] at (fm.north) {erase};
\node[font=\scriptsize, above=3pt] at (ad.north) {write};
\node[draw, rounded corners, fill=white, minimum height=7mm] (bus) at (0,-0.4) {$h_{t-1},\ x_t$};
\draw[gray] (bus.east) -- (9,-0.4);
\node[gt] (fg) at (3,0.6) {forget\\$f_t$};
\node[gt] (ig) at (5,0.6) {input\\$i_t$};
\node[draw, rounded corners, fill=lightnyupurple!30, minimum width=11mm, minimum height=7mm, align=center, inner sep=2pt] (cg) at (7,0.6) {cand.\\$\tilde{c}_t$};
\node[op] (ml) at (6,1.9) {$\times$};
\draw[gray] (3,-0.4) -- (fg.south);
\draw[gray] (5,-0.4) -- (ig.south);
\draw[gray] (7,-0.4) -- (cg.south);
\draw[->] (fg.north) -- (fm.south);
\draw[->] (ig.north) -- (5,1.9) -- (ml.west);
\draw[->] (cg.north) -- (7,1.9) -- (ml.east);
\draw[->] (ml.north) -- (ad.south);
\fi
\ifnum#1>2
\node[gt] (og) at (9,0.6) {output\\$o_t$};
\draw[gray] (9,-0.4) -- (og.south);
\fill (8.2,3) circle (2pt);
\node[op, minimum size=8mm] (th) at (8.2,1.9) {$\tanh$};
\node[op] (mo) at (9,1.9) {$\times$};
\node[st] (ht) at (11,1.9) {$h_t$};
\node[font=\scriptsize, gray, below] at (ht.south) {output};
\draw[->] (8.2,3) -- (th.north);
\draw[->] (th.east) -- (mo.west);
\draw[->] (og.north) -- (mo.south);
\draw[->] (mo.east) -- (ht.west);
\fi
\end{tikzpicture}}}

%%%%%%%%%%%%%%%%%%%%%%%%%%%%%%%%%%%%%%%%%%%%%%%%%%%
\begin{frame}
  \begin{center}
    {\Large Long Short-Term Memory (LSTM)}
  \end{center}
\end{frame}


%%%%%%%%%%%%%%%%%%%%%%%%%%%%%%%%%%%%%%%%%%%%%%%%%%%
\begin{frame}[fragile] \frametitle{What an LSTM Adds}
\begin{itemize}
\item In a plain RNN or a GRU, one vector $h_t$ does two jobs: it is the memory and it is the output
\item The LSTM (Hochreiter and Schmidhuber, 1997) separates them
  \begin{itemize}
  \item \textbf{Cell state} $c_t$: the long-term memory, a conveyor belt that runs through time
  \item \textbf{Hidden state} $h_t$: what the cell shows at this step
  \end{itemize}
\item Three gates, each set on its own: \textbf{forget} (erase), \textbf{input} (write) and \textbf{output} (show)
\item The problem it targets: carrying information across a long gap, for example the number of load cycles since the last service, without every step's work overwriting it
\end{itemize}
\begin{block}{Intuition}
Think of a notebook and what you say aloud. The notebook (cell state) keeps what matters over a long time; what you say at each step (hidden state) is only the part of the notebook that is useful right now.
\end{block}
\end{frame}


%%%%%%%%%%%%%%%%%%%%%%%%%%%%%%%%%%%%%%%%%%%%%%%%%%%
\begin{frame}[fragile] \frametitle{The Cell-State Highway}
\begin{center}
\lstmdiagram{1}
\end{center}
\begin{itemize}
\item The memory $c_t$ runs along the top, from $c_{t-1}$ to $c_t$
\item It is changed at only two points: scaled down (erase), then added to (write)
\item No weight matrix and no $\tanh$ sits on this path, so information can ride it nearly unchanged
\end{itemize}
\end{frame}


%%%%%%%%%%%%%%%%%%%%%%%%%%%%%%%%%%%%%%%%%%%%%%%%%%%
\begin{frame}[fragile] \frametitle{Forget Gate, Input Gate and Candidate}
\begin{center}
\lstmdiagram{2}
\end{center}
\begin{itemize}
\item All of them read the same pair $(h_{t-1}, x_t)$, drawn once as the line at the bottom
\item $f_t = \sigma(W_{fx}x_t + W_{fh}h_{t-1} + b_f)$ is how much of the old memory to keep; $i_t = \sigma(W_{ix}x_t + W_{ih}h_{t-1} + b_i)$ is how much of the candidate to write
\item $\tilde{c}_t = \tanh(W_{cx}x_t + W_{ch}h_{t-1} + b_c)$ and $c_t = f_t \odot c_{t-1} + i_t \odot \tilde{c}_t$
\item One slot: $c_{t-1} = 2.0$, $f_t = 0.9$, $i_t = 0.5$, candidate $0.6$ gives $c_t = 0.9 \cdot 2.0 + 0.5 \cdot 0.6 = 2.1$
\end{itemize}
\end{frame}


%%%%%%%%%%%%%%%%%%%%%%%%%%%%%%%%%%%%%%%%%%%%%%%%%%%
\begin{frame}[fragile] \frametitle{The Output Gate}
\begin{center}
\lstmdiagram{3}
\end{center}
\begin{itemize}
\item $o_t = \sigma(W_{ox}x_t + W_{oh}h_{t-1} + b_o)$ and $h_t = o_t \odot \tanh(c_t)$
\item $\tanh$ squeezes the memory into $(-1, 1)$; the output gate decides how much of it to show
\item The cell can store a value now and reveal it only when it is needed
\item Same slot: $\tanh(2.1) = 0.9705$ and $o_t = 0.5$ give $h_t \approx 0.49$
\item Both $c_t$ and $h_t$ go to the next step; $h_t$ is also what the next layer sees
\end{itemize}
\end{frame}


%%%%%%%%%%%%%%%%%%%%%%%%%%%%%%%%%%%%%%%%%%%%%%%%%%%
\begin{frame}[fragile] \frametitle{LSTM and Vanishing Gradients}
\begin{itemize}
\item Along the cell-state path, the factor from one step to the previous one is the forget gate: $\partial c_t / \partial c_{t-1} = f_t$ (slot by slot)
\item In a plain RNN the factor involves $W_{hh}$ and the slope of $\tanh$, usually below 1: in the earlier test the first input's effect fell to about $10^{-21}$ after 80 steps
\item If the network learns $f_t \approx 0.99$, then over 50 steps about 61\% of the gradient survives ($0.99^{50} = 0.61$)
  \begin{itemize}
  \item With $f_t = 0.9$ only 0.5\% survives ($0.9^{50} = 0.005$): the gate must learn to stay near 1
  \end{itemize}
\item Gates ease the problem; they do not remove it
  \begin{itemize}
  \item At the untrained start a GRU and an LSTM also fade (about $10^{-17}$ at 80 steps, same test as the plain RNN)
  \item Gates give the network a way to hold information once trained, not a free pass at the start
  \end{itemize}
\end{itemize}

{\tiny (Ref: How LSTM networks solve the problem of vanishing gradients
- Nir Arbel)}

\end{frame}


%%%%%%%%%%%%%%%%%%%%%%%%%%%%%%%%%%%%%%%%%%%%%%%%%%%
\begin{frame}[fragile] \frametitle{What an LSTM Gives and What It Costs}
\textbf{Gives}
\begin{itemize}
\item Long-range memory held apart from the output
\item Separate control to erase, write and show
\item Well tested on many sequence tasks
\end{itemize}
\textbf{Costs}
\begin{itemize}
\item Four blocks of weights: about four times a plain RNN, against three for a GRU
\item Two states to carry, more computation, more to tune
\item Still one step at a time, and very long gaps remain hard
\item On small problems a GRU often does as well
\end{itemize}
\end{frame}


%%%%%%%%%%%%%%%%%%%%%%%%%%%%%%%%%%%%%%%%%%%%%%%%%%%
\begin{frame}[fragile] \frametitle{Plain RNN, GRU and LSTM Side by Side}
\begin{center}
\small
\adjustbox{max width=\linewidth}{%
\begin{tabular}{lccc}
\hline
 & Plain RNN & GRU & LSTM \\
\hline
States carried & 1 & 1 & 2 ($h_t$, $c_t$) \\
Gates & none & 2 & 3 \\
Weights, against a plain RNN & 1x & 3x & 4x \\
Keep and replace & no choice & tied, one gate & separate gates \\
Memory apart from output & no & no & yes \\
\hline
\end{tabular}}
\end{center}
\end{frame}


%%%%%%%%%%%%%%%%%%%%%%%%%%%%%%%%%%%%%%%%%%%%%%%%%%%
\begin{frame}[fragile]\frametitle{Quick Check: Gated Cells}
\begin{block}{Think About It}
1. In a GRU, the update gate $z_t$ is close to 1 for one memory slot. What happens to that slot?
\begin{enumerate}
  \item[A.] It is replaced by the candidate.
  \item[B.] It keeps its old value almost unchanged.
  \item[C.] It is set to zero.
  \item[D.] It is shown only to the next layer.
\end{enumerate}
2. What does an LSTM have that a GRU does not?
\begin{enumerate}
  \item[A.] Gates.
  \item[B.] A separate cell state, so memory is kept apart from the output.
  \item[C.] Immunity to vanishing gradients at the untrained start.
  \item[D.] Fewer parameters.
\end{enumerate}
\end{block}
\end{frame}

%%%%%%%%%%%%%%%%%%%%%%%%%%%%%%%%%%%%%%%%%%%%%%%%%%%
\begin{frame}[fragile]\frametitle{Quick Check: Gated Cells}
\begin{block}{Answer}
\textbf{1. Correct: B.} The new state is $z_t h_{t-1} + (1 - z_t)\tilde{h}_t$, so $z_t$ near 1 keeps the old value. A describes $z_t$ near 0.

\textbf{2. Correct: B.} A GRU has gates too (A), and at the untrained start both cells fade (C). An LSTM has more parameters, not fewer (D); its distinct feature is the separate cell state $c_t$.
\end{block}
\end{frame}
